\documentclass[11pt]{article}
\usepackage[a4paper,margin=25mm]{geometry}
\usepackage{amsmath,amssymb,amsthm}
\usepackage[hidelinks]{hyperref}
\newtheorem{theorem}{Theorem}
\newtheorem{lemma}[theorem]{Lemma}
\newtheorem{corollary}[theorem]{Corollary}
\theoremstyle{remark}
\newtheorem*{remark}{Remark}
\newcommand{\Sur}{\mathrm{Sur}}
\title{Surjective Homomorphisms from the Plactic Monoid:\\
the Optimal Exponent Is $n$.\\
A Proof of Conjecture 00000003841}
\author{GitHub: gaochengzhecpu}
\date{}
\begin{document}
\maketitle
\begin{abstract}
Let $P_n$ be the plactic monoid on $n\ge1$ letters. We prove that a real
number $c$ has the property that every finite monoid $M$ receives at
most $|M|^c$ surjective homomorphisms from $P_n$ if and only if
$c\ge n$. So a constant $c(n)$ as in the source exists, the optimal one
is $c(n)=n$, it is at least $1$, and it is unbounded in $n$. The
statement and its proof are formalised in Lean with Mathlib, including
the plactic monoid and real exponents.
\end{abstract}

\section*{Statement and reading of the source}
Let $A_n=\{1<2<\dots<n\}$ and let $A_n^*$ be the free monoid of words.
Knuth equivalence is the monoid congruence generated by
\[
 acb\equiv cab\ \ (a\le b<c),\qquad bac\equiv bca\ \ (a<b\le c),
\]
and the plactic monoid is $P_n=A_n^*/{\equiv}$. For a monoid $M$ let
$\Sur(P_n,M)$ be the set of surjective monoid homomorphisms
$P_n\to M$. Call a real number $c$ \emph{admissible for $n$} if
\[
 |\Sur(P_n,M)|\le|M|^{c}\qquad\text{for every finite monoid }M .
\]
The source, in both languages, asserts three things: an admissible
constant $c(n)$ exists; the optimal exponent satisfies $c(n)\ge1$; and
the optimal exponent grows unboundedly with $n$. We read ``optimal
exponent'' as the least admissible real number. The source does not
say that $c(n)$ is an integer, so real exponents are allowed.

\begin{theorem}\label{thm:main}
Let $n\ge1$. A real number $c$ is admissible for $n$ if and only if
$c\ge n$.
\end{theorem}

\begin{corollary}\label{cor:conj}
For every $n\ge1$ the least admissible exponent exists and equals $n$.
In particular $c(n)\ge1$, and $c(n)\to\infty$. The conjecture holds.
\end{corollary}

\section*{Proof}
\begin{lemma}\label{lem:upper}
For every $n\ge0$ and every finite monoid $M$, the number of all
homomorphisms $P_n\to M$ is at most $|M|^n$.
\end{lemma}
\begin{proof}
The classes $[1],\dots,[n]$ of the letters generate $P_n$, because
every word is a product of letters. So a homomorphism $f\colon P_n\to
M$ is determined by the tuple $(f[1],\dots,f[n])\in M^n$.
\end{proof}

\begin{lemma}\label{lem:comm}
Let $M$ be a commutative monoid. Every map $g\colon A_n\to M$ extends
to a homomorphism $\bar g\colon P_n\to M$ with $\bar g[i]=g(i)$.
\end{lemma}
\begin{proof}
The map $g$ extends to a homomorphism $A_n^*\to M$. In a commutative
monoid the two sides of each Knuth relation have the same image, since
they are products of the same three elements. So the extension is
constant on Knuth classes.
\end{proof}

\begin{lemma}\label{lem:lower}
Let $p$ be a prime and $n\ge0$. Then
$|\Sur(P_n,\mathbb Z/p)|\ge p^n-1$, where $\mathbb Z/p$ is the cyclic
group of order $p$ regarded as a monoid.
\end{lemma}
\begin{proof}
Let $g\colon A_n\to\mathbb Z/p$ be a map that is not identically $0$,
say $g(i)=a\ne0$. Let $b\in\mathbb Z/p$, and let $k\in\{0,\dots,p-1\}$
represent $ba^{-1}$. Then $\bar g([i]^k)=ka=b$. So $\bar g$ is
surjective. Different maps $g$ give different homomorphisms, and there
are $p^n-1$ maps that are not identically $0$. (In fact equality holds,
since the homomorphism with $g=0$ is not surjective.)
\end{proof}

\begin{proof}[Proof of Theorem~\ref{thm:main}]
Let $c\ge n$ and let $M$ be a finite monoid. Then $|M|\ge1$, so by
Lemma~\ref{lem:upper}
$|\Sur(P_n,M)|\le|M|^n\le|M|^c$. So $c$ is admissible.

Let $c<n$ and suppose that $c$ is admissible. Put $\delta=n-c>0$ and
choose a prime $p\ge3$ with $p\ge2^{1/\delta}$. Then $p^\delta\ge2$,
so $p^c=p^n/p^\delta\le p^n/2$. Admissibility for $M=\mathbb Z/p$ and
Lemma~\ref{lem:lower} give
\[
 p^n-1\le|\Sur(P_n,\mathbb Z/p)|\le p^c\le p^n/2 ,
\]
hence $p^n\le2$. But $p^n\ge p\ge3$ because $n\ge1$. This contradiction
shows that $c$ is not admissible.
\end{proof}

\begin{proof}[Proof of Corollary~\ref{cor:conj}]
By Theorem~\ref{thm:main} the set of admissible exponents for $n\ge1$
is $[n,\infty)$. Its least element is $n\ge1$, and for every real bound
$B$ all admissible exponents for all $n\ge\max(B,1)$ are at least $B$.
\end{proof}

\begin{remark}
(1) \emph{The case $n=0$.} $P_0$ is the trivial monoid. It has one
surjective homomorphism to a trivial monoid and none to any other, so
every real number is admissible and there is no least one. The source
speaks of the alphabet $\{1,\dots,n\}$, and the statement about the
optimal exponent is proved for $n\ge1$.

(2) \emph{What is used about $P_n$.} The proof uses two properties: $P_n$
is generated by $n$ elements, and the Knuth relations hold in every
commutative monoid. The same proof therefore gives the same answer for
the free monoid, the free commutative monoid and the hypoplactic monoid
on $n$ letters. The result is elementary; the fine structure of the
plactic monoid does not enter.

(3) \emph{Counting quotients.} The label of the source is ``plactic
quotient count bound'', while its statement in both languages counts
surjective homomorphisms, and this is what is proved here. If one
counted surjective homomorphisms up to automorphisms of $M$, that is,
congruences on $P_n$ with quotient isomorphic to $M$, then $|M|^n$ is
still a bound and $\mathbb Z/p$ has $(p^n-1)/(p-1)\ge p^{n-1}$ such
congruences, so the optimal exponent would lie in $[n-1,n]$ and would
still be unbounded; but for $n=1$ it would be $0$, so the clause
$c(n)\ge1$ would hold only from $n=2$ on. This variant is not
formalised.
\end{remark}

\section*{Formal verification and scope}
\texttt{KnuthRel} is the pair of Knuth relations on Mathlib's
\texttt{FreeMonoid} of a linearly ordered alphabet, \texttt{placticCon}
is the congruence generated by it (Mathlib's \texttt{conGen}), and
\texttt{Plactic n} is the quotient monoid for the alphabet
\texttt{Fin n}. Homomorphisms are Mathlib's monoid homomorphisms, and
\begin{center}\small
\texttt{numSurj n M}
\end{center}
is the cardinality of the type of surjective homomorphisms from
\texttt{Plactic n} to \texttt{M}. Lemma~\ref{lem:upper} is
\texttt{hom\_ext\_gen} and \texttt{numSurj\_le}. Lemma~\ref{lem:comm}
is the definition \texttt{liftComm} with \texttt{liftComm\_gen}.
Lemma~\ref{lem:lower} is
\begin{center}\small
\texttt{surjective\_liftComm},\quad \texttt{pow\_le\_numSurj\_add\_one},
\end{center}
for the group \texttt{Multiplicative (ZMod p)}.

\texttt{IsAdmissible n c} is the displayed inequality for all finite
monoids, with $c$ a real number and Mathlib's real power.
Theorem~\ref{thm:main} is \texttt{isAdmissible\_iff}, and the
least-element statement is \texttt{isLeast\_admissible}. The proposition
\texttt{Conjecture} is the conjunction of: for every $n$ there is an
admissible exponent; for every $n\ge1$ there is a least admissible
exponent and it is at least $1$; for every real $B$ there is $N$ such
that for all $n\ge N$ every admissible exponent is at least $B$. The
final theorem is
\begin{center}\small\texttt{conjecture\_true : Conjecture}.\end{center}
No \texttt{sorry}, \texttt{native\_decide} or custom axiom occurs.

Everything in Theorem~\ref{thm:main} and Corollary~\ref{cor:conj} is
formalised. The quantification is over monoids in the lowest universe,
which contains a copy of every finite monoid. The three points of the
remark are not formalised.

\section*{Source and reproduction}
The exact bilingual source is included byte-for-byte as
\texttt{SOURCE.md}. Its repository path is
\begin{center}\small\texttt{conjectures/00000003841.md}.\end{center}
The project pins Lean 4.19.0 and Mathlib commit
\begin{center}\small\texttt{c44e0c8ee63ca166450922a373c7409c5d26b00b}.\end{center}
From \texttt{lean/}, run \texttt{lake build}, then
\begin{center}\small\texttt{lake env lean -DwarningAsError=true Main.lean}.\end{center}
The included public-Git manifest locks all dependency revisions. The
\texttt{verification/} directory records the fresh-build logs,
theorem-axiom reports, artifact hashes, review and final PDF
inspection. No supplementary numerical computation is required.
\end{document}
